% Build tables: python Reports/M3/build_tables.py (from repository root).
% Build PDF: xelatex M3_writing_brief.tex twice in Reports/M3.
\documentclass[11pt,a4paper]{article}
\usepackage[margin=2.1cm]{geometry}
\usepackage{fontspec}
\setmainfont{Times New Roman}
\usepackage{amsmath,amssymb,bm}
\usepackage{booktabs,tabularx,array}
\usepackage{enumitem,xcolor}
\usepackage[hidelinks]{hyperref}

\newcolumntype{Y}{>{\raggedright\arraybackslash}X}
\renewcommand{\arraystretch}{1.18}
\setlist{itemsep=3pt,topsep=4pt}
\setlength{\parindent}{0pt}
\setlength{\parskip}{5pt}
\setlength{\emergencystretch}{2em}
\setlength{\headheight}{14pt}
\pagestyle{plain}
\newcommand{\E}{\mathbb E}
\newcommand{\Prb}{\mathbb P}
\newcommand{\note}[1]{\par\smallskip\textcolor{teal!70!black}{\textit{Writing note:} #1}\par}
\newcommand{\file}[1]{{\small\texttt{\detokenize{#1}}}}
% Generated by Reports/M3/build_tables.py. Do not edit numeric tables.

\newcommand{\InitialTable}{%
\begin{tabular}{@{}lrrrrr@{}}
\toprule
Cohort & Area (ha) & Planting year & $V_0$ (\%) & $M_0$ (mm) & Quota (mm/yr) \\
\midrule
A & 100 & 1 & 17.114 & 14.265 & 300 \\
B & 50 & 1 & 23.446 & 20.106 & 600 \\
C & 150 & 2 & 16.903 & 14.492 & 300 \\
\bottomrule
\end{tabular}}

\newcommand{\BaseTable}{%
\begin{tabular}{@{}rrrrrr@{}}
\toprule
Year & First hit (\%) & 95\% Wilson interval & Low area (ha) & Ever-hit area (ha) & $H_{\rm eco}$ \\
\midrule
10 & 77.47 & [76.392, 78.506] & 106.07 & 213.12 & 1168.68 \\
30 & 99.57 & [99.366, 99.704] & 155.62 & 292.72 & 3093.88 \\
50 & 100.00 & [99.936, 100.000] & 200.22 & 299.45 & 4761.89 \\
\bottomrule
\end{tabular}}

\newcommand{\CohortTable}{%
\begin{tabular}{@{}lrrrrr@{}}
\toprule
Entity & $p_{50}$ (\%) & RMST (yr) & MC SE (yr) & Censored & Final $V$ (\%) \\
\midrule
A & 100.00 & 7.182 & 0.077 & 0 & 9.043 \\
B & 98.90 & 16.107 & 0.147 & 66 & 13.768 \\
C & 100.00 & 7.131 & 0.077 & 0 & 9.043 \\
Project & 100.00 & 7.131 & 0.077 & 0 & 9.831 \\
\bottomrule
\end{tabular}}

\newcommand{\SensitivityTable}{%
\begin{tabular}{@{}lrrrr@{}}
\toprule
Scenario & $p_{10}$ (\%) & $p_{50}$ (\%) & RMST (yr) & $H_{\rm eco}(50)$ \\
\midrule
Base & 77.47 & 100.00 & 7.13 & 4761.9 \\
Stationary & 75.62 & 99.95 & 7.54 & 5360.7 \\
No storms & 0.00 & 0.00 & 50.00 & 7106.7 \\
No drought & 73.83 & 100.00 & 7.86 & 5070.9 \\
No random disasters & 0.00 & 0.00 & 50.00 & 7271.8 \\
Stable, no disasters & 0.00 & 0.00 & 50.00 & 7439.5 \\
No irrigation & 100.00 & 100.00 & 2.96 & 1026.2 \\
Water request 80\% & 81.00 & 100.00 & 6.54 & 4299.1 \\
Water request 120\% & 72.30 & 100.00 & 8.08 & 5125.5 \\
Start-year shocks & 65.27 & 100.00 & 9.45 & 4732.8 \\
Independent storms & 81.20 & 100.00 & 6.56 & 4778.1 \\
No drought--storm link & 75.10 & 100.00 & 7.64 & 4924.6 \\
Fixed loss per event & 73.60 & 100.00 & 8.12 & 4831.1 \\
Threshold $0.20K$ & 35.07 & 98.70 & 17.08 & 4795.4 \\
Threshold $0.40K$ & 92.50 & 100.00 & 4.40 & 4795.4 \\
Wind trend only & 76.17 & 100.00 & 7.69 & 5250.4 \\
Rain trend only & 76.07 & 99.93 & 7.79 & 5371.4 \\
Wetter, calmer & 73.20 & 99.83 & 8.20 & 5633.3 \\
\bottomrule
\end{tabular}}

\newcommand{\JointTable}{%
\begin{tabular}{@{}rrrr@{}}
\toprule
$\lambda_0$ & $\mathrm{E}[U]$ & $p_{50}$ (\%) & RMST (yr) \\
\midrule
0.10 & 0.15 & 44.87 & 41.78 \\
0.10 & 0.30 & 97.90 & 16.65 \\
0.10 & 0.45 & 99.73 & 11.38 \\
0.25 & 0.15 & 97.73 & 20.68 \\
0.25 & 0.30 & 100.00 & 7.27 \\
0.25 & 0.45 & 100.00 & 5.17 \\
0.50 & 0.15 & 100.00 & 8.78 \\
0.50 & 0.30 & 100.00 & 3.74 \\
0.50 & 0.45 & 100.00 & 2.82 \\
\bottomrule
\end{tabular}}

\newcommand{\ChainTable}{%
\begin{tabular}{@{}lrrr@{}}
\toprule
Cohort & Stationary mean (yr) & Variance (yr$^2$) & RMST(50) (yr) \\
\midrule
A & 7.852 & 49.323 & 7.846 \\
B & 21.747 & 398.477 & 19.971 \\
C & 7.782 & 49.093 & 7.776 \\
\bottomrule
\end{tabular}}

\newcommand{\StateError}{0.0380}

\newcommand{\ShockError}{0.2024}

\newcommand{\MCError}{1.055}

\newcommand{\Threshold}{11.423}

\begin{document}
\thispagestyle{plain}
\begin{center}
{\LARGE\bfseries Model 3 Writing Brief}\par
{\large Post-construction resilience under compound disasters}\par
{\large Results, model logic and handoff to Claude}\par
2026-11-13
\end{center}

This brief follows the M1/M2 report format. It is a writing and plotting guide, not a new model run.
All numeric result tables are extracted from \file{Results/v7/model3/}; the extraction script verifies
the recorded input and source-code hashes before writing them. No M1/M2 result has been changed.
No scientific figure is produced here: Claude owns the figure design and rendering.

\section{Goal and completed work}
\begin{tabularx}{\textwidth}{@{}p{3.4cm}YY@{}}
\toprule Task 3 need & Implementation & Delivered evidence \\\midrule
Build on Tasks 1 and 2 & Carry each M2 cohort's final cover and storage into the unchanged M1 water--growth map & Three cohorts; interface replay error below $10^{-5}$ \\
Random disasters & Full compound-Poisson wind/sand loss, conditional on shared drought; drought changes rain, demand and supply & Joint event simulations; no one-event approximation \\
Changing background & Wind-to-demand response, rainfall trend, event-rate trend and drought-probability trend & Nonstationary baseline and stationary controls \\
Long-term performance & First passage, current low-cover area and cumulative cover integral & 50-year tables; 27 scenarios \\
Numerical verification & Daily water balance, resource caps, stationary 2-D chain, Monte Carlo and grid checks & 28 M1/M2/M3 tests passed; separate error checks \\
Writing and figures & Equations, assumptions, numerical results and exact figure data mapping & This brief plus the Chinese handoff in this folder \\
\bottomrule
\end{tabularx}

\paragraph{Main finding.}
Under the declared baseline assumptions, the probability that at least half the project area has
simultaneously crossed below the diagnostic threshold at some annual census is about 77.47\% by
post-year 10 and 99.57\% by post-year 30. All 6,000 simulated paths trigger this event within 50 years.
Yet the expected area currently below threshold at year 50 is only 200.22 ha, whereas 299.45 ha has
ever crossed it. Recovery makes these different quantities. The 50-year cover integral is
4,761.89 effective ha$\cdot$yr, versus 7,271.79 without random disasters under the same background trends.

\paragraph{Interpretation.}
These are \textbf{conditional stress-test results}, not calibrated local forecasts.
The event parameters, future trends and degradation threshold are assumptions. The numerical work is
complete, but \file{formal_model_ready=false} remains appropriate. A first crossing is neither permanent
destruction nor the physical end of an engineering asset.

\paragraph{Reading route.}
Sections 2--4 explain the state, equations and algorithms; Sections 5--7 provide parameters,
results and validation; Sections 8--9 specify figures and defensible writing claims.

\clearpage
\section{Construction and interfaces}
\subsection*{Clock and initial state}
Post-year 0 is the end of M2 project year 15. Post-year 1 is project year 16, not calendar year 2027.
Use the exported final states, not $V=K$ or a reset storage bucket.
\begin{center}\small\InitialTable\end{center}
The cohort aliases used throughout this report are:
\begin{itemize}
\item A: \file{eco_y1_100_to_300_after_3}, 100 ha.
\item B: \file{eco_y1_ramp_100_300_600}, 50 ha.
\item C: \file{eco_y2_100_to_300_after_3}, 150 ha.
\end{itemize}
These quotas are applications for water, not guaranteed deliveries. Cohort area stays fixed.
The inherited 9.5 MW PV installation requests 190 m$^3$/yr for cleaning. Continuing O\&M and cleaning
are scenario assumptions; no 50-year PV revenue or replacement-cost forecast is claimed.

\subsection*{One physical year, in execution order}
\begin{enumerate}
\item Apply the deterministic background trends for that post-construction year.
\item Draw one regional drought state; it changes precipitation, evapotranspiration demand,
water-supply caps and the conditional storm-event rate.
\item Draw the full Poisson event count and all event-loss marks. Baseline cohorts share the storm
realization. The independent-cohort sensitivity still shares drought.
\item Reserve PV cleaning, then ration ecological requests against monthly and annual supply.
Run the M1 daily bucket and yearly vegetation growth for every cohort.
\item Apply the compound storm retention at year end. Record annual states, first crossing and
current low-cover area. Keep evolving the physical state even after a first crossing.
\end{enumerate}
The start-year-shock sensitivity moves storm loss before the M1 update; the failure census is
still at the end of the complete year. Annual observation can miss a short-lived within-year crossing.

\subsection*{What changed from the old M3 draft}
The old kernel used $1-\lambda$ and $\lambda$ while claiming a compound Poisson process.
It omitted storage memory, initialized cover at $K$, and mixed row and column transition conventions.
Those issues are corrected. The average retention formula is exact under iid marks, not a small-loss
approximation. Zero storm rate does not imply zero failure if irrigation stops. The historical,
unidentified M1 wind-loss coefficient is not restored or amplified.

\note{A structure figure should show M2 terminal states $\to$ joint drought/storm inputs $\to$
monthly allocation and M1 daily bucket $\to$ annual census $\to$ risk/benefit summaries.
Draw a separate validation branch to the stationary 2-D chain. Do not label the MC physical state as absorbing.}

\clearpage
\section{Equations and their meaning}
\subsection*{Disasters and background}
For year index $t=0,\ldots,49$, let $D_t$ be the shared drought indicator:
\begin{align}
D_t&\sim\operatorname{Bernoulli}(p_{D,0}+b_Dt), &
N_t\mid D_t&\sim\operatorname{Poisson}(\lambda_0 e^{\alpha_\lambda t}\gamma^{D_t}),\\
U_{tq}&\sim\operatorname{Uniform}(u_{\rm lo},u_{\rm hi}), &
R_t&=\prod_{q=1}^{N_t}(1-U_{tq}),\\
\E[R_t\mid D_t]&=\exp\{-\lambda_0e^{\alpha_\lambda t}\gamma^{D_t}\E[U]\}.&&
\end{align}
With no event, $R_t=1$. Wind abrasion, sandstorm damage and burial are aggregated into this loss;
they are not counted as three independent hits. The last equation alone identifies only a product
of frequency and mean loss, not their separate values. Their full distribution can carry more information.

Wind baseline scales by $e^{\alpha_ut}$ and rain by $e^{\alpha_Rt}$.
The existing M1 Penman--Monteith response ratio converts the wind change into monthly demand.
Drought additionally multiplies rain by 0.40, demand by 1.15 and supply caps by 0.70.
This is a synthetic climate sequence, not a downscaled climate-model forecast.

\subsection*{Water and vegetation}
Let $\Phi$ denote the M1 annual map, including daily water updates:
\begin{align}
(\widetilde V_{j,t+1},M_{j,t+1})&=\Phi(V_{j,t},M_{j,t};P_t,ET_{0g,t},I_{j,t,m}),\\
\widetilde V_{j,t+1}&=V_{j,t}+rV_{j,t}(1-V_{j,t}/K)\overline W_{j,t}-dV_{j,t},\\
V_{j,t+1}&=R_t\widetilde V_{j,t+1},\\
M_{\rm end}&=M_{\rm start}+P_{\rm eff}+I_{\rm net}+G-T_a-D_{\rm drain}.
\end{align}
Here $\overline W$ is actual/potential transpiration, not an external wetness trend.
Gross irrigation is converted using $\eta_w$ before entering storage. One mm over one ha is 10 m$^3$.
Supply limits apply to all maintained cohorts plus PV cleaning, including cohorts that previously crossed the threshold.

\subsection*{Risk, recovery and ecological integral}
\begin{align}
\tau_j&=\inf\{t\ge0:V_{j,t}<\rho K\},\\
\tau_{\rm proj}&=\inf\left\{t:\frac{\sum_j a_j\mathbf1[V_{j,t}<\rho K]}{\sum_j a_j}\ge\beta\right\},\\
\operatorname{RMST}(H)&=\E[\min(\tau,H)]=\sum_{t=0}^{H-1}\Prb(\tau>t),\\
H_{\rm eco}(H)&=\sum_{t=0}^{H-1}\sum_j a_jV_{j,t}/K.
\end{align}
Baseline $\rho=0.30$, $\beta=0.50$. Thus the threshold is \Threshold\% cover proxy, not 30\% absolute cover.
The ecological integral uses start-of-year cover, as in M2; the start-shock sensitivity uses cover after
that initial shock. It is a cover proxy, not measured erosion prevention. Initial failure has $\tau=0$.
Non-failures through year 50 are right-censored; code 51 in the histogram is not an observed failure year.

\clearpage
\section{Algorithms and reproducibility}
\subsection*{A. Direct simulation: the nonstationary main result}
\begin{enumerate}
\item Load the final M1 parameters and M2 cohorts. Replay every M2 cohort through current M1;
stop if a cover/storage discrepancy exceeds $10^{-5}$.
\item For each scenario, initialize all paths with the same three inherited states and use seed 20261113.
Use 6,000 paths for baseline/stationary cases, 3,000 for each other case.
\item Each year draw drought and the full compound storm retention. Calculate wet/dry monthly allocation
ledgers. Update only the paths belonging to each drought state with the vectorized daily M1 map.
\item At the annual census, update first-passage flags and physical states separately.
Accumulate ecological integrals and export means, quantiles, event frequencies and sample paths.
\item Estimate probabilities with Wilson 95\% intervals. Export RMST and its MC standard error,
plus the explicit number of paths censored at the horizon.
\end{enumerate}
Common seeds across cases aid reproducibility but do not guarantee perfectly paired shocks after
changing the draw count or array shape. The case differences are descriptive sensitivity results,
not a paired hypothesis test.

\subsection*{B. Stationary killed chain: a numerical cross-check}
Discretize $(V,M)$ over $[\rho K,K]\times[0,C]$. Apply the same water allocation and M1 map, then test
failure in continuous state before bilinear projection of surviving states. Average over drought
and the full compound-loss distribution to form the substochastic matrix $Q$.
\begin{align}
\Prb(\tau\le H)&=1-\bm p_0Q^H\bm1,\\
(I-Q)\bm t&=\bm1, & (I-Q)\bm s&=2\bm t-\bm1,\\
\E[\tau]&=\bm p_0\bm t, & \operatorname{Var}(\tau)&=\bm p_0\bm s-(\bm p_0\bm t)^2.
\end{align}
$\bm p_0$ is a row distribution. The moment equations require finite mean absorption time;
they are not applied indiscriminately to zero-hazard cases. The no-event atom is exact. Positive-event
retention uses log-loss convolution and quantile compression; a Poisson tail below $10^{-12}$ is
assigned zero retention. The main Monte Carlo does not truncate the count.

\subsection*{C. Entry points}
\begin{tabularx}{\textwidth}{@{}p{4.9cm}Y@{}}
\toprule File / function & Purpose \\\midrule
\file{Code/v7_m3_core.py} & \file{load_inputs}, \file{batch_year}, \file{allocation}, \file{storm_draw}, \file{simulate}, \file{stationary_chain} \\
\file{Code/v7_model3.py} & Full 27-case export, three chain grids, provenance hashes \\
\file{Code/v7_m3_grid_check.py} & Cross-grid check separating state and quadrature errors \\
\file{Code/test_v7_m3.py} & Nine M3 tests; 28 when combined with M1/M2 \\
\file{Reports/M3/build_tables.py} & Verify provenance and extract this report's tables only \\
\bottomrule
\end{tabularx}
From repository root: \file{python Code/v7_model3.py}, then
\file{python Code/v7_m3_grid_check.py}. Dependencies are in \file{Code/requirements-m3.txt}.
Plotting needs neither command if hashes still match. Rebuild only this brief's tables with
\file{python Reports/M3/build_tables.py}; compile this folder's TeX twice with XeLaTeX.

\clearpage
\section{Parameters, constraints and evidence}
\begin{tabularx}{\textwidth}{@{}p{3.3cm}p{4.0cm}Y@{}}
\toprule Quantity & Baseline & Evidence / interpretation \\\midrule
Growth $r$, attrition $d$ & 0.546698; 0.107306 & Inherited M1 fit; uncertainty not propagated here \\
$K$; storage capacity $C$ & 38.0754\%; 75 mm & M1 regional proxy and effective bucket assumption \\
Water efficiency; rain fraction & 0.75; 0.25 & M1 assumptions; groundwater baseline 0 \\
Initial states; area; quotas & Table in Section 2 & M2 output, not new measured stands \\
Annual water cap & $1.5\times10^6$ m$^3$ & Inherited M2 resource scenario; drought factor applies \\
Monthly caps, Jan--Dec & 50, 60, 90, 120, 180, 220, 240, 220, 150, 90, 50, 30 (thousand m$^3$) & Inherited M2 scenario; enforced monthly and annually \\
Storm rate $\lambda_0$ & 0.25 yr$^{-1}$ & Assumed frequency; not derived from monthly mean wind \\
Event loss $U$ & Uniform [0.10, 0.50] & Assumed fractional cover loss; not a measured local damage distribution \\
Drought probability & $0.12+0.002t$ & 0.12--0.218 over simulated years; no serial persistence \\
Drought modifiers & Rain 0.40; ET demand 1.15; supply 0.70 & Assumed joint stress; not fitted drought response \\
Drought storm-rate multiplier & $\gamma=1.50$ & Assumed event association, not a correlation coefficient \\
Trend coefficients & Wind +0.005; rain $-0.005$; storm rate +0.010 per year & Log-rate scenario changes; time starts at 0 \\
Diagnostic thresholds & $\rho=0.30$; $\beta=0.50$ & Management diagnostics, not verified protection standards \\
Horizon and seed & 50 years; 20261113 & Numerical choices, not fixed by the problem \\
State / shock grids & (81,11,24), (161,21,48), (321,41,96) & Counts of cover nodes, storage nodes and positive-event quantile nodes \\
Cross-grid and tail & (161,21,96); $10^{-12}$ & Isolate approximation errors \\
\bottomrule
\end{tabularx}

\paragraph{Hard checks in the implementation.}
Probabilities must remain in $[0,1]$ for every simulated year; invalid trends are rejected rather than
silently clipped. Loss marks satisfy $0\le u_{\rm lo}\le u_{\rm hi}<1$; event rates are nonnegative and
may exceed one event per year. States remain in $0\le V\le K$, $0\le M\le C$.
Rationing preserves the monthly and annual caps; a 20\% higher request is not a 20\% higher guaranteed delivery.

\paragraph{Evidence boundary.}
Historical forcing and the M1 response are reused. FAO-56 supports the wind dependence of reference
evapotranspiration, not the selected future trends or local damage rates. The event process needs a
dated disaster catalogue and before/after cover observations for calibration. M1 background attrition
may already contain some historical disaster loss, so overlap with the extra stress process remains unresolved.
The full field-by-field register is \file{Data/v7/m3_inputs/README.md}.

\clearpage
\section{Results and interpretation}
\subsection*{Baseline joint project event}
\begin{center}\small\BaseTable\end{center}
The first-hit column is cumulative probability. Low area is expected area currently below threshold;
ever-hit area counts a cohort if it has crossed at any previous annual census.
$H_{\rm eco}$ is in effective ha$\cdot$yr. All rows refer to the 6,000-path baseline.
The 100\% sample hit rate at year 50 still has a nonzero-width Wilson interval.

\subsection*{Cohort-level results at post-year 50}
\begin{center}\small\CohortTable\end{center}
The 50 ha cohort B retains higher mean cover and a longer restricted first-passage time under the
baseline. This does not establish an optimal adaptation policy: areas and quotas are inherited from M2,
not reoptimized in M3. The project event uses simultaneously low area, not a product of cohort probabilities.
Its coincidence with cohort C's baseline summary must not be assumed for other configurations.

\subsection*{Why first passage and recovery both matter}
\begin{itemize}
\item The first-passage flag stays true after a crossing; the physical vegetation state can recover.
\item High cumulative crossing probability does not mean that the same area remains degraded at every later date.
\item RMST(50) is an average of $\min(\tau,50)$ across all paths, including censored ones.
It is not the mean among failures only, nor the unconditional physical service life of the project.
\end{itemize}

\subsection*{Stationary reference, finest grid}
\begin{center}\small\ChainTable\end{center}
These means and variances belong only to the stationary reference with continuing hazard and maintenance.
They are from the 321-by-41 state grid and 96 positive-event quadrature nodes. They are not the
nonstationary baseline's infinite-horizon life expectancy. Section 7 compares chain probabilities with
continuous-state Monte Carlo and separates state-grid and shock-integration errors.

\note{The paper should report a horizon, an event definition and the evidence class with every risk number.
For example: ``Under the declared stress scenario, the 10-year probability of a first annual census
with at least half of treated area below $0.30K$ is 77.47\%.''}

\clearpage
\subsection*{Mechanism, resource and threshold sensitivities}
\begin{center}\small\SensitivityTable\end{center}
Stationary means all four trends are zero. No random disasters retains the baseline climate trends;
stable/no-disasters removes both. Wind-only and rain-only retain stationary random disasters.
Water cases multiply applications before resource rationing. Threshold cases alter the diagnostic,
not the physical vegetation map. Fixed-loss uses $U=0.30$ per event, keeping the base mean but changing
the mark distribution. These contrasts are scenario comparisons, not causal estimates from field data.

\subsection*{Joint frequency--severity experiment}
\begin{center}\small\JointTable\end{center}
Each cell uses uniform marks on $[0.5\mu,1.5\mu]$, where $\mu=\E[U]$; other baseline settings remain.
The $\lambda_0=0.25,\mu=0.30$ cell therefore has range [0.15,0.45], unlike the baseline [0.10,0.50].
It must not be presented as an exact duplicate of the baseline. Only nine cells are simulated;
do not turn them into a finely resolved empirical surface. Show RMST or 10-year risk when 50-year
probabilities saturate. At the low-rate/low-loss corner, the project 50-year first-hit probability
is about 44.87\%, illustrating strong dependence on uncalibrated assumptions.

\clearpage
\section{Verification and limits of the checks}
\begin{tabularx}{\textwidth}{@{}p{4.2cm}Y@{}}
\toprule Check & Observed evidence \\\midrule
Regression tests & 28 passed: 12 M1, 7 M2, 9 M3. M3 covers scalar/vector M1 agreement, allocation, Poisson moments, initial failure, censoring, invalid input and chain comparison. \\
M1/M2 interface replay & Maximum cover/storage discrepancy $2.824\times10^{-6}$, below $10^{-5}$ tolerance. M2 tables were preserved. \\
Monthly water mass balance & Maximum residual below $4.84\times10^{-13}$ mm. \\
Resource and export checks & 27 scenarios, 108 entity summaries, 32,400 conditional monthly rows; all caps satisfied within rounding tolerance. Histograms reproduce counts, censoring and RMST. \\
Finest chain vs. MC & All 9 cohort-by-horizon points inside MC 95\% intervals; maximum absolute probability difference \MCError\ percentage points. \\
State-grid refinement & Holding 96 shock nodes fixed, 161$\times$21 to 321$\times$41: maximum probability change \StateError\ percentage points. \\
Shock quadrature refinement & Holding 161$\times$21 fixed, 48 to 96 nodes: maximum change \ShockError\ percentage points. \\
Provenance & 12 input/source hashes checked; report tables extracted without simulation. \\
\bottomrule
\end{tabularx}

The medium grid has one comparison point outside the MC interval; it is retained in the export.
The refinement differences are below the one-percentage-point numerical criterion, but do not prove
convergence for every possible parameter set. MC intervals themselves have sampling variability.
Do not write ``less than one percentage point error against MC'': the largest observed difference is
\MCError\ percentage points.

\subsection*{Three uncertainty objects that must remain separate}
\begin{enumerate}
\item \textbf{Path spread:} 5th--95th percentiles describe random outcomes conditional on the assumed parameters.
They are not a 95\% confidence interval for the mean.
\item \textbf{Simulation error:} Wilson probability intervals and RMST standard errors reflect finite MC sample size.
They exclude parameter and structural uncertainty.
\item \textbf{Model uncertainty:} M1 estimation, cover proxies, initial states, event rates/marks, climate trends,
threshold choice, spatial dependence and serial drought dependence are not jointly propagated here.
\end{enumerate}

\paragraph{Controls are not proofs of field accuracy.}
Zero storms with continued water can remain above threshold in these runs; zero storms with stopped
water can fail. The latter is an explicit test against the old ``$\lambda=0$ means infinite lifetime'' rule.
Grid agreement and water closure establish numerical consistency, not empirical local validity.

\clearpage
\section{Figure handoff to Claude}
\textbf{Status: data ready; figures below are proposed and await rendering.}\par
Read only \file{Results/v7/model3/} after checking the saved hashes. Use the existing paper palette and
MATLAB conventions if suitable; layout and final figure selection are Claude's responsibility.

\begingroup\small
\begin{tabularx}{\textwidth}{@{}p{0.6cm}p{3.7cm}Y@{}}
\toprule ID & Intended message & Exact data and plotting contract \\\midrule
F1 & Model structure (schematic) & Section 2 execution order; show M2 states, shared drought/storm, monthly rationing, M1 bucket, annual census and the separate stationary-validation branch. No invented spatial layout. \\
F2 & Cohort cover paths & \file{cohort_annual.csv}; base case, three cohorts, \file{V_mean_pct}, \file{V_p05_pct}, \file{V_p95_pct}; threshold \Threshold\%. Label the band as path percentiles. \\
F3 & First passage versus current condition & Same table: \file{failure_probability} versus \file{current_below_threshold_probability}; or two area columns in \file{project_annual.csv}. Do not confuse an area expectation with a probability. \\
F4 & Background and cumulative benefit & \file{project_annual.csv}; base, stationary, no-random-disasters. Plot project first-hit probability and \file{cumulative_cover_ha_year_mean}; optional p05/p95 benefit band. \\
F5 & Joint hazard sensitivity & \file{summary.csv}; \file{entity=project}, nine \file{joint_rate_*} cases. Read axes from \file{scenario_configs.json}; paired grids for $p_{50}$ and RMST, or 10-year risk from annual data. \\
F6 & Water and event assumptions & Selected Section 6 sensitivity cases; \file{summary.csv} and annual project tables. Keep threshold changes distinct from physical management changes. \\
F7 & Supply rationing and climate & \file{resources.csv}; request versus delivery, gross/cap ratio and annual $\sum P/\sum ET_{0g}$. Wet/dry rows are alternative conditional ledgers, not additive water use. \\
F8 & Numerical validation & \file{markov_mc_comparison.csv}, \file{grid_validation.csv}; distinguish state refinement, quadrature refinement and MC intervals. Only stationary results. \\
\bottomrule
\end{tabularx}
\endgroup

\paragraph{Optional path illustration.}
\file{sample_paths.csv} contains only the first 12 paths per case. Use it to explain a recovery trajectory,
with a declared selection rule; never recompute risk from this illustrative subset.

\paragraph{Caption template.}
``Conditional post-construction scenario, initialized at M2 year 15 end. The diagnostic threshold is
$0.30K$ (\Threshold\% cover proxy). Baseline/stationary simulations use 6,000 paths; sensitivity cases use
3,000. Bands represent [path percentiles / Wilson MC intervals, as appropriate]. Hazard parameters
and future trends are assumed rather than locally calibrated.''

\paragraph{Delivery check.}
Keep post-year and project-year labels explicit; preserve percentage versus fraction units; area-weight
project means. Do not plot censor code 51 as a failure year or turn the 100\% hit rate into a zero-width
interval. No long-term NPV, carbon, sediment tonnes or actual protected-area claims are supported by these tables.

\clearpage
\section{Writing claims, paper outline and handoff files}
\small
\begin{tabularx}{\textwidth}{@{}YY@{}}
\toprule Defensible statement & Statement to avoid \\\midrule
The M2 plan was evaluated under specified compound hazards and background trends. & We predict the actual disaster history of this region over the next 50 years. \\
The first-crossing risk is high under the baseline assumptions; physical recovery remains possible. & The project will certainly be permanently destroyed. \\
RMST is the mean truncated first-passage time in a 50-year assessment window. & RMST is the engineering service life or a calibrated expected lifetime. \\
The high-water cohort has higher cover and longer RMST in this scenario. & M3 proves that this irrigation allocation is economically optimal. \\
Numerical cross-checks agree to the reported tolerances. & Agreement with Monte Carlo validates the model against real disasters. \\
The ecological integral is inherited from M2's cover proxy. & The integral measures tonnes of erosion prevented or verified protected land. \\
\bottomrule
\end{tabularx}

\subsection*{Suggested paper narrative}
\begin{enumerate}
\item \textbf{Question and interface.} Why a deterministic construction plan needs a resilience stress test;
inherit the three terminal M2 states and declare the post-construction clock.
\item \textbf{Stochastic extension.} Introduce shared drought and compound wind/sand loss, resource rationing
and background trends, using a short structure figure and the main update equations.
\item \textbf{Metrics.} Define cohort and project first passage, current low-cover area, RMST and cover integral.
Explain right censoring and recovery before quoting numerical risk.
\item \textbf{Main results.} Present baseline paths, risk and ecological integral, with a stationary and
no-random-disaster comparison. State the assumptions beside the results.
\item \textbf{Sensitivity and reliability.} Present joint hazard and selected water/ordering/dependence
contrasts, then a compact numerical-validation panel. Avoid nine similar low-information figures.
\item \textbf{Implications and limitations.} Recommend monitoring and evidence collection based on sensitive
inputs; do not claim an optimized adaptive repair policy that was never solved.
\end{enumerate}

\subsection*{What evidence would improve the next version}
A dated, geolocated disaster catalogue; before/after cover and survival observations; local supply
and maintenance commitments; field evidence for a protection threshold; and a defensible climate
projection would allow calibration rather than assumption-driven stress testing. Repeated drought,
partial spatial correlation, subannual failures and funded repair decisions are extensions, not completed features.

\subsection*{Repository handoff}
\begin{itemize}
\item \file{Reports/M3/README.md}: Chinese reading guide, headline findings and current figure handoff.
\item \file{Reports/M3/M3_writing_brief.tex}: editable report; \file{m3_result_tables.tex}: generated numeric tables.
\item \file{Docs/20261113_m3_audit_and_implementation.md}: correction history and numerical audit.
\item \file{Docs/20261113_m3_figure_handoff.md}: detailed field-level Chinese plotting instructions.
\item \file{Reports/model_v7/m3_correction.tex}: current formal modeling chapter; this brief complements it with results.
\end{itemize}
\end{document}
